\section{Bluesync: a Time Synchronization protocol for BLE}
Paper \cite{asgarian2021bluesync} explains an interesting protocol that can be used for time synchronization purposes when dealing with Bluetooth Low Energy devices, like sensor network nodes. The article stresses the importance of a low-cost and high-efficiency mechanism that can work on these resource-constrained IoT nodes.  

In particular, BlueSync is defined as a synchronization service for BLE beacon devices. BLE Beaconing is an advertising mode, supported by Bluetooth Low Energy, that enables the exchange of data between devices without the need of entering a bonded connection: the server device will \textit{advertise} its data, while the client can read it doing a \textit{scanning} operation. This protocol does not present any time synchronization service. 
BlueSync makes the system work in two different stages: first, there is a timeslot (few seconds) in which synchronization packages are exchanged, then, the second timeslot will allow the execution of synchronous tasks. Of course, the time interval between the two stages should be decided with care, as precision will depend on that. The developers showed that the node can proceed with 10 minutes of synchronous tasks without the need of re-synchronization, still with a good time precision. 

The devices tested in the work are from Nordic Semiconductor, in particular an nRF51 (16~MHz ARM Cortex-M0) and an nRF52 (64~MHz ARM Cortex-M4).

\paragraph{Implementation} The implementation of the BlueSync protocol is presented in three steps, that are reported and explained as follows.
\begin{itemize}
    \item Timestamping. The main problem encountered when dealing with BLE beacons is that they usually exchange data in three different channels, with non-predictable delays. With proper modifications of the registers of the Nordic devices, the developers managed to set the transmission (advertising) and the reception (scanning) of the data on a single channel, both on the server and the client devices. Once this has been resolved, the problem of delays is addressed and solved by precisely deciding the moment at which packets are timestamped. Also in this case, a particular function available on Nordic devices has been used: PPI, Programmable Peripheral Interconnect. It allows the execution of a particular task at the end of a transmission event in a very precise manner, generating a standard deviation in timestamping of less than 50~ns. Of course, the timestamps collected after the radio event are sent in the next packet, hence the synchronization stage requires more than one packet to be transmitted (to retrieve all the necessary timestamps).
    \item Frequency Drift Estimation. As soon as all synchronization packets have been received, the contained timestamps are used to estimate the offset and the drift of the nodes clocks. This is a calculation that is performed on the node itself, and should be optimized due to their hardware constraints and available functionalities (for example, the Floating Point Unit could be missing). Researchers tried two different algorithms: Linear Regression and Average Error. Both have been tested with either 8 or 16 timestamps, sent at regular intervals, with the latest packet initiating the execution of the synchronous task. The two algorithms show different performances on the two different processors, due to the different architectures: when dealing with heterogeneous hardware, it is important to fine-tune the choice of the algorithm. 
    \item Tick Adjustment. The last part of the mechanism consists in applying the estimation calculated before to the number of ticks necessary to reach the next point in time to execute the synchronous task: the floating-point result should be converted to an integer and then set as timer interrupt. These adjustment require clock cycles, and their impact should be taken into consideration, for example recalculating the ticks every 100 runs.
\end{itemize}

\paragraph{Conclusions} To sum it up, BlueSync is a synchronization protocol aimed at BLE Sensor Networks, compatible with commercial SoCs. Its implementation consists mainly in improving timestamping of BLE packets and choosing the most optimized frequency-drift estimation algorithm. The resulting timing accuracies, in the tests performed in the laboratory, are around 320~ns for 60~s. They assure that the system can provide good accuracy with synchronization timeslots every 10 minutes, with synchronization windows of less than 2 seconds. Even if these results are particularly good, it is important to notice that they are possible only with the functionalities available on the particular Nordic devices used, hence they may not be applicable to other SoCs on the market. 